\scanpage{300}
\mspage{216}

\margin{Pour ces réductions, cf 4.15.8 (p.453 à \unsure{468|463}). L'ultime réduction se trouve dans \unsure{4.15.8.9} (p.~\unsure{460|468}), et son corollaire 4.15.10.}% slanted

J'ai pu me ramener, par des arguments similaires,$\vphantom{x}^{\ast}$ au cas où le monoïde $M=\operatorname{End}\unsure{(\xi)|(\zeta)}$ est un monoïde de monomorphismes i.e.

\[
uv=uv'\Longrightarrow v=v'\qquad\text{si }u,v,v'\in M,
\]

ce \deleted{\illegible{1}} qui implique que les flèches de transition dans le système inductif des $M_{\unsure{i|ki}}=M$ sont des mono. Il s'ensuit que l'on ne peut avoir $\operatorname{card}F(\xi)\leq\pi_0$ ($I$ étant grand devant $\pi_0$) que si $\alpha_{j i_0}$ est un iso pour \inserted{$j\geq i_0$, et} $i_0$ assez grand, i.e. si le syst. inductif \deleted{est \unsure{ess. co}} induit \deleted{par} pour $\deleted{\unsure{i}}j\geq i_0$ est constant, a fortiori que le syst. inductif envisagé est \unsure{ess.t|esst.} constant. Mais je viens de passer près de quatre heures pour essayer de prouver ce point, sans y être arrivé ! Cela semble plus délicat à prouver que je ne le prévoyais, mais je crois toujours que c'est bel et bien vrai.

\margin{$\unsure{\xi|\rbrace}$}
% archive sheet number: 299
